% Section 5.7, from lectures/L05/appendix/b_exercises.md.

\section{Exercises}
\label{c5:sec:exercises}

\begin{aside}
\textbf{How to check your work.} Every exercise below that asks you to write a VHDL module has a
self-checking testbench under \file{lectures/L05/exercises/}. Write your module in its directory
\file{exercises/<module>/}, with the entity name and the \textbf{port order} the exercise states,
and then run it with GHDL \textemdash{} see \secref{c2:sec:testbenches} for the three commands.
\end{aside}

% ------------------------------------------------------------------------------------------
\exerciseset{\code{variable} versus \code{signal}}

\exercise{Trace the accumulator bug}{Reasoning}
\label{c5:ex:trace}
The broken \code{parity_gen} architecture from \secref{c5:sec:variable} is reproduced below:

\begin{vhdlcode}
architecture broken of parity_gen is
signal parity_s: std_logic;
begin
    parity <= parity_s;

    process(bits) is
    begin
        parity_s <= '0';
        for i in bits'range loop
            parity_s <= parity_s xor bits(i);
        end loop;
    end process;
end architecture;
\end{vhdlcode}

Assume \code{bits} has the range \code{7 downto 0}, that \code{bits = "10110110"}, and that
\code{parity_s} holds \code{'0'} before the process starts.

\xpart{a} Trace the process by hand. Note, for the initial assignment and for every loop iteration:
the current index \code{i}, the value of \code{bits(i)}, the value read from \code{parity_s}, the
value scheduled for \code{parity_s}, and whether that scheduled assignment survives until the
process suspends.

Remember that a signal assignment does not update the signal immediately, and that assignments to
the same signal during one and the same round replace earlier assignments scheduled for the same
simulation time.

\xpart{b} Determine the final value assigned to \code{parity_s} after the process suspends. Then
compute the correct XOR parity of \code{"10110110"} by hand, compare it with the result from the
broken architecture, and explain why repeated assignments to the signal do not create an
accumulator.

\xpart{c} Trace the variable-based architecture \code{behaviour} from \file{parity_gen.vhd} with the
same input. Note, for every loop iteration: the value of \code{bits(i)}, the value of \code{acc}
before the XOR operation, and the value of \code{acc} immediately after it. Confirm that \code{acc}
keeps each iteration's result so that the next iteration can use it, and that the final value
assigned to \code{parity} is the correct XOR parity.

\exercise{\code{min_of_two}}{VHDL}
\label{c5:ex:minoftwo}
Write an entity named \code{min_of_two}.

\begin{booktable}{@{}p{3.6em}p{3.4em}p{11.4em}L@{}}
  \thead{Port} & \thead{Dir.} & \thead{Type} & \thead{Meaning} \\
  \midrule
  \code{a} & in & \code{std_logic_vector(3 downto 0)} & First value. \\
  \code{b} & in & \code{std_logic_vector(3 downto 0)} & Second value. \\
  \code{m} & out & \code{std_logic_vector(3 downto 0)} & The smaller of the two. \\
\end{booktable}

Treat \code{a} and \code{b} as unsigned four-bit values. Implement the circuit with a process:
\begin{itemize}
  \item Include \code{a} and \code{b} in the sensitivity list.
  \item Declare a process-local variable holding the smaller value.
  \item Compare \code{a} and \code{b} as unsigned numbers, using the facilities in
    \code{ieee.numeric_std} per \secref{c2:sec:vectors}.
  \item Assign the variable exactly once per round, and assign it to \code{m} after the comparison.
\end{itemize}

This exercise needs no loop. Its point is to use a variable as temporary process-local storage for
something \emph{other} than an accumulator: \secref{c5:sec:variable}'s worked example builds a value
up across the iterations of a loop, which is the case a variable is most often introduced for, and
it is easy to leave thinking that is the only thing they are good for. Here the variable simply
holds an intermediate result for the few lines separating the computation from its use, which is the
more common of the two uses in real code.

\modulefig[0.62]{lectures/L05/appendix/images/min_of_two.png}

\begin{selfcheck}
\textbf{Self-check.} Name your entity \code{min_of_two}, with the ports \code{a}, \code{b} (in) and
\code{m} (out), all \code{std_logic_vector(3 downto 0)}, declared in that order; its testbench is in
\file{exercises/min_of_two/}.
\end{selfcheck}

\exercise{\code{ones_count8}}{VHDL}
\label{c5:ex:ones}
Write an entity named \code{ones_count8} that counts how many bits in an 8-bit input are set.

Exercise~\ref{c5:ex:trace} had you trace the accumulator bug on paper, and
Exercise~\ref{c5:ex:minoftwo} used a variable that was no accumulator at all. This one is the case
\secref{c5:sec:variable} is actually about: a value built up across the iterations of a loop. It is
the only exercise in the book where you write one.

\begin{booktable}{@{}p{4em}p{3.4em}p{11.4em}L@{}}
  \thead{Port} & \thead{Dir.} & \thead{Type} & \thead{Meaning} \\
  \midrule
  \code{bits} & in & \code{std_logic_vector(7 downto 0)} & The value to count over. \\
  \code{ones} & out & \code{natural range 0 to 8} & The number of bits in \code{bits} that are
    \code{'1'}. \\
\end{booktable}

Implement it with a single combinational process:
\begin{itemize}
  \item Put \code{bits} in the sensitivity list.
  \item Declare a process-local \code{variable} for the running count, and set it to \code{0} at
    the top of every round. A variable keeps its value between rounds, so a count that is never
    cleared keeps growing.
  \item Loop over \code{bits'range} and add \code{1} for every bit that is set.
  \item Assign the variable to \code{ones} once, after the loop.
\end{itemize}

Then answer, in one sentence each:
\begin{itemize}
  \item Write the same thing with a \code{signal} as the accumulator instead, run the testbench, and
    note what it reports for the input \code{00000010}. Explain the number you get using the
    scheduling rule, not by guessing.
  \item The loop runs eight times. How many adders does this design put on the FPGA, and how many
    clock cycles does it take to produce an answer? If either answer surprises you, reread
    \secref{c5:sec:variable}.
  \item \code{ones} is a \code{natural range 0 to 8} rather than a
    \code{std_logic_vector(3 downto 0)}. Both hold the answer. What does the range buy you, and what
    would GHDL do if your count ever left it?
\end{itemize}

\modulefig[0.66]{lectures/L05/appendix/images/ones_count8.png}

\begin{selfcheck}
\textbf{Self-check.} Name your entity \code{ones_count8}, with the input \code{bits}
(\code{std_logic_vector(7 downto 0)}) and the output \code{ones} (\code{natural range 0 to 8}),
declared in that order; its testbench is in \file{exercises/ones_count8/} and sweeps all 256 input
values.
\end{selfcheck}

% ------------------------------------------------------------------------------------------
\exerciseset{A module worth having on its own}

\exercise{\code{sync}: the synchroniser as a module of its own}{VHDL}
\label{c5:ex:sync}
Write the two-flip-flop synchroniser as a module of its own, named \code{sync}.

You have already built this circuit. It is the first two flip-flops of \chapref{c4:ch}'s
\code{button_sync}, written inline there because at that point it was the \emph{idea} that mattered:
an asynchronous input needs two flip-flops before anything else gets to look at it. Here you factor
those two out into a module, and the exercise is not the logic, which you already know, but
everything around it.

Two things make it worth doing again. The first is that a bare synchroniser is what most
asynchronous inputs actually need: \code{button_sync} adds an edge detector because a button press
is an event, but a serial line, a status flag from another clock domain, or a switch whose level is
all you ever read, need the two flip-flops and nothing more. \Chapref{c8:ch}'s serial receiver
instantiates exactly this on its \code{rx} pin. The second is that it takes \textbf{two generics},
and one of them is not a number.

\begin{booktable}{@{}p{4.4em}p{9.8em}p{4em}L@{}}
  \thead{Generic} & \thead{Type} & \thead{Default} & \thead{Meaning} \\
  \midrule
  \code{SIZE} & \code{natural range 1 to 15} & \code{1} & How many independent signals to
    synchronise, and therefore the width of both vector ports. \\
  \code{PRESET} & \code{std_logic} & \code{'0'} & The value both stages take at reset. \\
\end{booktable}

and these ports:

\begin{booktable}{@{}p{5.6em}p{3.4em}p{11em}L@{}}
  \thead{Port} & \thead{Dir.} & \thead{Type} & \thead{Meaning} \\
  \midrule
  \code{clock} & in & \code{std_logic} & 50 MHz system clock. \\
  \code{reset_s2_n} & in & \code{std_logic} & Active-low, \textbf{already synchronised} reset. A
    synchroniser does not synchronise its own reset; \chapref{c4:ch}'s \code{reset_sync} does
    that. \\
  \code{async_in} & in & \code{std_logic_vector(SIZE - 1 downto 0)} & The asynchronous input or
    inputs, straight from outside the clock domain. \\
  \code{sync_out} & out & \code{std_logic_vector(SIZE - 1 downto 0)} & The same signals, safe to
    use, two rising edges later. Every bit is synchronised independently. \\
\end{booktable}

\xpart{a} Implement it:
\begin{itemize}
  \item Two internal signals of type \code{std_logic_vector(SIZE - 1 downto 0)}, one per stage.
  \item One process, sensitive to \code{clock} and \code{reset_s2_n}.
  \item At reset, drive both stages to \code{PRESET}, outside the clocked branch.
  \item At a rising clock edge, shift \code{async_in} through the two stages.
  \item Drive \code{sync_out} from the second stage with a concurrent assignment.
\end{itemize}

\xpart{b} \code{PRESET} is a generic of type \code{std_logic}, the first in this book that is not a
number, and it exists because the safe reset value depends on what the signal \emph{means}. Only the
instantiating design knows that. Answer:
\begin{itemize}
  \item \code{button_sync} synchronises active-low buttons, so it would pass \code{'1'} in. What
    would \code{'0'} claim about the buttons during the two cycles after every reset?
  \item A serial line idles high, so \chapref{c8:ch}'s receiver also passes \code{'1'} in. What
    would \code{'0'} look like to a receiver watching for a start bit?
  \item Why is the \emph{default} \code{'0'} rather than \code{'1'} all the same, given both of
    those, and what does that say about whose job it is to get it right?
\end{itemize}

\xpart{c} This is also a question about what a generic costs. \code{SIZE} is
\code{natural range 1 to 15} and \code{PRESET} is \code{std_logic}, and both are fixed before
synthesis runs. Use \secref{c5:sec:fabric}'s picture of what the FPGA actually consists of and say
how many flip-flops an instance with \code{SIZE = 4} uses, how many an instance with
\code{SIZE = 1} uses, and where in the built circuit \code{PRESET} ends up. Neither of the two turns
up as logic anywhere; say what happened to them instead.

\xpart{d} Verify it with its testbench. It instantiates your module twice, once one bit wide with
\code{PRESET = '1'} and once three bits wide with \code{PRESET = '0'}, so a design that hardcodes
either generic fails. It checks that \code{sync_out} follows \code{async_in} after \textbf{exactly}
two rising edges rather than after one or eventually, that every bit in a vector is synchronised
separately, and that reset loads \code{PRESET} with no clock edge involved.

\note[Hint]Write the architecture for \code{SIZE = 1} in your head first, then check that every line
is already correct for a vector. The aggregate \code{(others => PRESET)} and an ordinary vector
assignment both work element by element, so the generalisation should cost no extra code.

\modulefig[0.76]{lectures/L05/appendix/images/sync.png}

\begin{selfcheck}
\textbf{Self-check.} Name your entity \code{sync}, with the generics \code{SIZE}
(\code{natural range 1 to 15}, default \code{1}) and \code{PRESET} (\code{std_logic}, default
\code{'0'}), the inputs \code{clock}, \code{reset_s2_n}, \code{async_in}
(\code{std_logic_vector(SIZE - 1 downto 0)}) and the output \code{sync_out}
(\code{std_logic_vector(SIZE - 1 downto 0)}), declared in that order; its testbench is in
\file{exercises/sync/}.

\textbf{Keep this file.} \Chapref{c8:ch}'s serial receiver capstone instantiates it on the \code{rx}
pin, and asks you to copy it in.
\end{selfcheck}

% ------------------------------------------------------------------------------------------
\exerciseset{From VHDL to hardware}

The Quartus flow is demonstrated on a real DE0-CV board during the sessions \textemdash{} first with
the brake assistant back in \chapref{c1:ch} \textemdash{} and is written down in
appendix~\ref{app:quartus} for reference. You need neither Quartus nor a board yourself; the
exercises below are about understanding what that flow does with the VHDL you write.

\exercise{What the toolchain does}{Reasoning}
\label{c5:ex:quartus}
Answer each of the following, with reference to appendix~\ref{app:quartus} and to what you saw
demonstrated. One or two sentences each.

\xpart{a} A Quartus project needs a chosen \textbf{top-level entity} before it compiles. What does
the tool do differently with the top-level entity than with any other entity in the project?
\code{min_of_two} from Exercise~\ref{c5:ex:minoftwo} and its testbench are both valid VHDL. Why can
only one of them ever be the top-level entity in an FPGA project?

\xpart{b} Pin assignment connects a port name to a physical pin. Why must the pins be assigned
\emph{before} compilation rather than afterwards? What would \code{m} physically be connected to if
you compiled \code{min_of_two} having assigned no pins at all?

\xpart{c} Synthesis turns your architecture into gates on the chip. \Chapref{c2:ch}'s exercise
\code{combo_logic} had you write $x = ab + c'd$ twice: once as a single expression, and once with
$c'd$ named by an internal signal. What should synthesis produce for each, and why is your answer
the same for both? Which of the two would you rather read in six months' time?

\xpart{d} Compilation gives warnings as well as errors. Why is ``it compiled without errors'' a much
weaker claim about an FPGA design than ``it compiled without errors'' is about a C program? Name one
kind of warning you would never ignore.

\exercise{Four predictions about the parity}{Simulation}
\label{c5:ex:parity}
\code{parity_gen} from \secref{c5:sec:variable} is the chapter's worked example, and it comes with a
reference testbench.

\xpart{a} Predict, before you run anything, \code{parity} by hand for each of these inputs:
\code{"00000000"}, \code{"10110110"}, \code{"11111111"} and \code{"10000000"}.

\xpart{b} Run the reference testbench against the worked example and confirm your four predictions:

\begin{shell}
cd lectures/L05/parity_gen
ghdl -a --std=93 parity_gen.vhd parity_gen_tb.vhd
ghdl -e --std=93 parity_gen_tb
ghdl -r --std=93 parity_gen_tb --assert-level=error --stop-time=10ms
\end{shell}

\xpart{c} Take any of those inputs and flip exactly one bit, leaving the other seven unchanged. What
happens to \code{parity}? Explain why changing exactly one input bit \emph{always} has to invert the
XOR parity result, whichever bit you pick and whatever the starting value. That property is what
makes a parity bit capable of detecting a single-bit error in transmission. What kind of error does
it miss, and why?

\exercise{A decoder that compiles cleanly and is still wrong}{Reasoning}
\label{c5:ex:latch}
A colleague sends you this combinational decoder for review. It compiles, and Quartus reports no
errors.

\begin{vhdlcode}
architecture behaviour of level_decode is
begin
    process(sel) is
    begin
        case sel is
            when "00"   => leds <= "0001";
            when "01"   => leds <= "0010";
            when "10"   => leds <= "0100";
            when others => null;
        end case;
    end process;
end architecture;
\end{vhdlcode}

\xpart{a} Quartus gives \code{Warning: Inferred latch(es) for signal "leds"}. Using
\secref{c5:sec:latch}, explain what the tool built and why. What did the VHDL ask for that a
combinational network cannot do?

\xpart{b} The author's defence is that \code{sel} is two bits, so \code{"11"} is the only value
\code{when others} can match, and the design never generates it. Give two separate reasons why that
argument does not hold: one about what \code{sel} can actually hold (\secref{c1:sec:vhdl}), and one
that would apply even if \code{sel} really could only be \code{"00"}, \code{"01"} or \code{"10"}.

\xpart{c} Fix it, in two different ways: by changing only the \code{when others} branch, and by
adding a single line \emph{before} the \code{case} statement and leaving every branch untouched.
Which of the two would you rather maintain, and why? The second scales to a \code{case} with twenty
branches; the first does not.

\xpart{d} This is a warning, not an error, and the design will often appear to work on the board.
Why does that make it more dangerous than an error rather than less?
